\begin{figure}[!htbp]
\begin{examplepanel}{srccol}{Source program \(P\) --- the running example}
\src{\(W=\top\to\top\). Let annotations specify the assigned types.}
\begin{lstlisting}[style=source]
let x         = { type A = W }
let f         = lambda (u: Top). u
let hidden    : x.A = f          -- L: W <: x.A
let recovered : W   = hidden     -- H: x.A <: W
recovered f                     -- T: W <: Top
\end{lstlisting}
\end{examplepanel}
\begin{examplepanel}{srccol}{Subtyping obligations}
\[
\src{
 \begin{array}{ccl}
 \text{L: hide} & W<:x.A & \text{the lower bound of }x.A,\\
 \text{H: recover} & x.A<:W & \text{the upper bound of }x.A,\\
 \text{T: pass the argument} & W<:\top & \text{inclusion into the universal type}.
 \end{array}}
\]
\end{examplepanel}
\caption{The running source program, with the object viewed at
\(\iface{W}{W}\). The aliases present the identity first at \(x.A\)
and then at \(W\); the final application presents its argument at
\(\top\). This sketch exposes the member bounds directly.
Section~\ref{sec:evidence} translates its client body;
Section~\ref{sec:closed-program} gives the complete translation, including
the recursive opening required by the formal source.}
\label{fig:program}
\end{figure}
